\chapter{Introduction}
\label{ch:introduction}

At its core, binary instrumentation is the technique of inserting additional code into an existing, compiled executable to observe or modify its runtime behavior without requiring access to the original source code. This capability is critical for enforcing security policies and monitoring performance.

\note{Advisor}{cyan}{Devo ancora estendere questo capitolo}
